% Part: history
% Chapter: biographies
% Section: ernst-zermelo

\documentclass[../../../include/open-logic-section]{subfiles}

\begin{document}

\olfileid{his}{bio}{zer}

\olsection{ایرنسٹ زرمَیلو}

\olphoto{zermelo-ernst}{ایرنسٹ زرمَیلو}

ایرنسٹ زرمَیلو 27 جولائی 1871 نوں جرمنی دے برلن شہر وچ جمیا۔ اوہدیاں پنج بھیناں سن، پر ٹبّر دی ماڑی صحت کرکے بس تِن وڈی عمر تک پہنچیاں۔ اوہدے ماں پیو وی اوہدی نِکّی عمر وچ فوت ہو گئے، تے جدوں اوہ ستاراں سالاں دا سی تاں اوہ تے اوہدے بھین بھرا یتیم ہو گئے۔ زرمَیلو نوں فنّاں وچ، خاص کر کے شاعری وچ، ڈونگھی دلچسپی سی۔ اوہ تِکھّی سمجھ، حاضر جوابی تے تنقید کرن لئی مشہور سی۔ اوہدیاں سبھ توں مشہور ریاضیاتی کامیابیاں وچ چُون دا مسلّمہ پیش کرنا (1904 وچ) تے سیٹ نظریے نوں مسلّمیاں نال بیان کرنا (1908 وچ) شامل نیں۔

یونیورسٹی وچ زرمَیلو دیاں دلچسپیاں کئی پاسے سن۔ اوہنے طبیعیات، ریاضی تے فلسفے دے مضمون پڑھے۔ ہرمان شوارٹس دی نگرانی تھلّے زرمَیلو نے اپنا مقالہ \emph{Investigations in
  the Calculus of Variations} 1894 وچ برلن یونیورسٹی توں پورا کیتا۔ 1897 وچ اوہنے گوٹنگن یونیورسٹی وچ ہور پڑھائی دا فیصلہ کیتا، جتھے ڈیوڈ ہلبرٹ دے ریاضی دیاں بنیاداں اُتے کم دا اوہدے اُتے بڑا اثر پیا۔ 1899 وچ اوہ پروفیسری لئی اہل ہو گیا، پر گیاراں سال پچھوں تک اوہنوں پروفیسری نہ ملی، شاید اوہدے اوپرے رویے تے «گھبراہٹ بھری جلدبازی» کرکے۔

آخر 1910 وچ زرمَیلو نوں زیورخ یونیورسٹی وچ تنخواہ والی پروفیسری ملی، پر تپِ دق کرکے 1916 وچ اوہنوں ریٹائر ہونا پیا۔ صحت ٹھیک ہون پچھوں 1921 وچ فرائی برگ یونیورسٹی وچ اوہنوں اعزازی پروفیسری دتی گئی۔ ایس سمیں اوہنے ریاضی دیاں بنیاداں اُتے کم کیتا۔ اوہ تھورالف سکولم تے کرٹ گوڈل دیاں لکھتاں توں کھِجھ گیا تے اپنے مقالیاں وچ اوہناں دے طریقیاں اُتے کھُلّھی تنقید کیتی۔ 1935 وچ اوہنوں فرائی برگ دے عہدے توں ہٹا دتا گیا، کیوں جے اوہ مقبول نہیں سی تے جرمنی وچ ہٹلر دے اقتدار وچ آؤن دا مخالف سی۔

زرمَیلو دی زندگی دے آخری سال تنہائی وچ لنگھے۔ 1935 وچ عہدے توں ہٹائے جان پچھوں اوہنے ریاضی چھڈ دتی۔ اوہ پینڈو علاقے وچ چلا گیا تے سادے ڈھنگ نال رہیا۔ اوہنے 1944 وچ ویاہ کیتا، تے نظر گھٹدی جان کرکے اوہ پوری طرح اپنی گھر والی اُتے ٹکن لگ پیا۔ 1951 تک زرمَیلو دی نظر پوری طرح چلی گئی۔ اوہ 21 مئی 1953 نوں جرمنی دے گُنٹرشٹال وچ وفات پا گیا۔

\begin{reading}
زرمَیلو دی پوری سوانح عمری لئی \citet{Ebbinghaus2015} ویکھو۔ زرمَیلو دے 1904 تے 1908 والے بنیادی مقالے اصل جرمن وچ پڑھے جا سکدے نیں \citep{Zermelo1904,Zermelo1908}۔ زرمَیلو دیاں اکٹھیاں لکھتاں، طبیعیات اُتے اوہدی لکھت سمیت، انگریزی ترجمے وچ \citep{Ebbinghaus2010,Ebbinghaus2013} اندر ملدیاں نیں۔
\end{reading}

\end{document}
